% Generated draft: requires mathematical and source review.
\documentclass[11pt]{article}
\usepackage[margin=25mm]{geometry}
\usepackage{fontspec}
\setmainfont{Latin Modern Roman}
\usepackage{amsmath,amssymb,mathtools,unicode-math}
\setmathfont{Latin Modern Math}
\usepackage{enumitem,needspace,xurl,hyperref}
\setlength{\parindent}{0pt}
\setlength{\parskip}{5pt}
\setlength{\emergencystretch}{3em}
\newcommand{\N}{\mathbb{N}}
\newcommand{\Z}{\mathbb{Z}}
\newcommand{\Q}{\mathbb{Q}}
\newcommand{\R}{\mathbb{R}}
\newcommand{\C}{\mathbb{C}}
\newcommand{\F}{\mathbb{F}}
\newcommand{\A}{\mathbb{A}}
\newcommand{\PP}{\mathbb{P}}
\newcommand{\E}{\mathbb{E}}
\newcommand{\eps}{\varepsilon}
\newcommand{\dd}{\,\mathrm{d}}
\begin{document}
\section*{UnsolvedMath ID: 700227}
\textbf{Source problem code:} catalogue entry
\subsection{Definitions and mathematical statement}
A \emph{variety of groups} is a class \(\mathfrak V\) consisting of all groups satisfying some set \(\Sigma\) of group identities. A group identity is an equation between group words, required to hold for every assignment of its variables to elements of the group. Neither the groups in \(\mathfrak V\) nor the set \(\Sigma\) are required to be finite. Define
\[
\widetilde{\mathfrak V}=\{G\mid G\text{ is a finite group and }G\in\mathfrak V\}.
\]
Classes of groups are understood up to isomorphism.

\textbf{Problem (R. Baer, 1973).} Characterize the classes \(\mathcal C\) of finite groups for which there exists a variety of groups \(\mathfrak V\) such that
\[
\mathcal C=\widetilde{\mathfrak V}.
\]
Equivalently, seek necessary and sufficient conditions on \(\mathcal C\) for the existence of a set \(\Sigma\) of group identities such that, for every finite group \(G\),
\[
G\in\mathcal C\quad\Longleftrightarrow\quad G\models\Sigma.
\]
The source does not prescribe a particular form for the requested characterization.

The supplied registry labels this problem open. Its background also reports an explicitly unverified AI claim of a characterization; that claim is not treated here as an established solution.
\subsection{Short English statement}
Characterize exactly which classes of finite groups occur as the finite members of a variety defined by group identities. The identities must hold under every substitution, and the variety may also contain infinite groups. The supplied open label and the reported unverified AI solution claim are retained without independent assessment.
\subsection{Sources}
\begin{itemize}
\item Local UnsolvedMath registry, permanent ID \(700227\), exact source problem code \texttt{catalogue entry}, title \emph{Kourovka Notebook Problem 4.17}; supplied complete source record.
\item R. Baer, Problem 4.17, \emph{The Kourovka Notebook (Unsolved Problems in Group Theory)}, 4th issue (1973). The supplied record cites the 21st issue edition (2026), \texttt{arXiv:1401.0300v48}. The registry identifies the Issues 1--21 collection with editors E. I. Khukhro and V. D. Mazurov.
\item Supplied discussion: Achyuth Jayadevan, \emph{Preprint} (2026), \url{https://zenodo.org/records/22724488}. The record explicitly describes this as an \emph{unverified AI solution claim} asserting that a characterization is given.
\item Supplied status review dated \texttt{2026-10-16}: canonical status \texttt{open}, evidence category \texttt{unrefereed\_claim}. Its notes report a \(\#\) marker for the unverified AI solution claim in the October 2026 21st issue, \texttt{arXiv:1401.0300v48}. This is supplied dated curation, not an independently verified literature assessment.
\end{itemize}
\end{document}
